%% Chapter 10, Section 10.3 — Discussion and Bibliography
%% Source: A First Course in Monte Carlo Methods, Sanz-Alonso & Al-Ghattas
%% Book pages 122–123

\section{Discussion and Bibliography}
\label{sec:10.3}

We refer to \cite{28,29} for accessible introductions to VI that have inspired the presentation in
this chapter. For a more in-depth exploration of VI, see \cite{119,233}. The paper \cite{112}
contains a modern approach to VI with applications to massive text data-sets. Applications to
graphical models, and a generalization of the mean-field family to allow for interactions, are
discussed in \cite{119}. While CAVI provides a simple framework for ELBO maximization under the
mean-field family, computing the CAVI updates can be challenging for highly complex targets, in
which case alternative optimization methods based on stochastic VI or autodifferentiation may be
required \cite{112,131}. Beyond the mean-field family, other simple but practical choices of
variational family include Gaussians \cite{208} and Gaussian mixtures \cite{133}.

Our presentation has solely focused on VI techniques that seek to minimize the KL-divergence
\cite{132}. The works \cite{141,110} consider VI using the broad family of Renyi-alpha divergences,
which includes for instance the KL-divergence, the $\chi^2$-divergence, and the Hellinger distance
\cite{86}. However, as emphasized in this chapter, the KL-divergence holds a special place in VI,
since among a wide class of divergences it is the unique choice for which minimization of the
objective does not require knowledge of the normalizing constant of the target \cite{47}. For
sampling problems, the paper \cite{79} explores gradient flows arising from minimizing various
objectives, and shows that the choice of KL-divergence also leads to fundamental practical
advantages.

An overarching idea behind VI, the EM algorithm, and other popular computational methods such as
variational auto-encoders \cite{123}, is to minimize KL-divergence by maximizing the ELBO\@. The
ELBO functional is typically non-convex, even for simple potentials; see for instance the discussion
in \cite{4}. Due to this lack of convexity, it is often challenging to develop theory for VI outside
restricted settings on the target distribution. For this reason, the study of statistical and
optimization convergence guarantees for VI is still an active area of research, see e.g.\ \cite{242}
for statistical convergence rates and \cite{133} for a theoretical framework that relies on gradient
flows.

The EM algorithm was introduced in \cite{55}, although particular instances had been previously
considered on numerous occasions, see e.g.\ \cite{43,106}. Convergence of the EM algorithm was
established in \cite{239}. We refer to \cite{158} for a textbook on the EM algorithm and to
\cite{27,107,152,28} for gentle introductions. EM is an example of a minimization-maximization
method \cite{115}, a broad family of algorithms with rich theoretical underpinnings. Several
modifications to the original EM algorithm have been proposed to speed up its convergence. We refer
to \cite{159,143,117,172} for some important methodological developments and to \cite{160} for a
survey written on the 20th anniversary of the original paper \cite{55}.
